%% kapa: exam=1 round=21 year=2010 period=1 form=A subject=economics subject_ko=경제원론 seq=30 number=70
%% source: Q-Net 감정평가사 기출문제 (cst00309) artlSeq=5024181 file=extracted/1008/5024181/2010년 제21회 감정평가사 1차시험 기출문제/1교시.hwp
%% conversion: kordoc 4.24.51 -> tools/gen1.py (auto); review=proofread-sonnet; flags=none
\begin{question}{70}
자연실업률에 관한 설명으로 옳은 것을 모두 고른 것은?
\begin{framedtext}
ㄱ. 자연실업은 구조적 실업과 경기적 실업의 합계를 말한다. \par
ㄴ. 자연실업률은 실제 실업률이 상승ㆍ하락하는 기준이 되는 정상적인 실업률이다. \par
ㄷ. 마찰적 실업의 증가는 자연실업률을 증가시킨다. \par
\end{framedtext}
\begin{choices}
  \item ㄱ
  \item ㄴ
  \item ㄷ
  \item ㄱ, ㄴ
  \item ㄴ, ㄷ
\end{choices}
\end{question}
\answer{5}
